% ============================================================
%  GUÍA 5: LÍMITES Y CONTINUIDAD
%  Curso de Introducción a la Universidad — Precálculo y Cálculo I
%  Autor: Gabriel Astudillo
%  Sesiones cubiertas: 13, 14, 15 y 16
% ============================================================

\documentclass[12pt, a4paper]{article}

% ── Paquetes ─────────────────────────────────────────────────

\usepackage{mathwizards-guia}
\mwguia{Guía 5 — Límites y Continuidad}{Gabriel Astudillo $\cdot$ Curso Preuniversitario}

\begin{document}
% ============================================================

% ── Portada / Título ─────────────────────────────────────────
\mwportada{Guía de Estudio 5}{Límites y Continuidad}{Límites Laterales, Cálculo de Límites, Límites Trigonométricos, Límites al Infinito y Continuidad}

\vspace{4pt}
\begin{cuadroinfo}
  \small
  \textbf{Presentación:} Los límites son el cimiento sobre el cual se construye todo el cálculo diferencial. Entender la idea de aproximación te permitirá luego comprender la derivada con profundidad. En esta guía trabajarás con límites desde una perspectiva intuitiva, luego algebraica, y finalmente con las técnicas para límites trigonométricos y al infinito. Cada sección incluye \textbf{ejercicios resueltos} y \textbf{propuestos}.
  \\[4pt]
  \textbf{Nivel de Dificultad:}\quad
  \facil{} Básico / Directo \quad
  \medio{} Intermedio \quad
  \dificil{} Desafiante
\end{cuadroinfo}

\vspace{8pt}
\begin{cuadroextra}
  \textbf{El porqué de los límites: la llave al mundo de lo infinitamente pequeño}\\[4pt]
  ¿Has notado que un velocímetro marca 90\,km/h justo en un instante? ¿Cómo se puede medir velocidad en un solo instante si la velocidad ``se calcula'' dividiendo distancia entre tiempo? La respuesta está en los \textbf{límites} y, con ellos, en los \textbf{infinitésimos}: cantidades que se acercan a cero sin ser exactamente cero, pero que capturan la esencia del movimiento y el cambio.

  \textbf{Dos ideas que lo cambiaron todo:}
  \begin{itemize}[leftmargin=*]
    \item \textbf{El polígono que sueña con ser círculo.} Hace más de dos mil años, Arquímedes imaginó polígonos regulares inscritos en un círculo: primero un triángulo, luego un hexágono, un dodecágono\dots Cuantos más lados agregaba, más se ``pegaba'' el polígono al círculo. Si pudieras repetir este proceso infinitamente, cada lado se volvería un \textbf{infinitésimo} de arco y el polígono se fundiría con el círculo. El límite nos permite decir sin paradojas: \textit{«el área del polígono \textbf{tiende} al área del círculo cuando el número de lados \textbf{tiende a infinito}»}. De esta idea nace el \textbf{cálculo integral}, que estudia cómo acumular piezas diminutas para obtener un todo.
    \item \textbf{La instantánea que esconde la velocidad.} Cuando calculas la velocidad media $v_m = \frac{\Delta x}{\Delta t}$ y tomas intervalos de tiempo $\Delta t$ cada vez más pequeños alrededor de un instante $t_0$, esa velocidad media se acerca a un único valor: justo lo que lees en el velocímetro. Eso es exactamente $\lim_{\Delta t \to 0} \frac{\Delta x}{\Delta t}$, el corazón de la \textbf{derivada} y del \textbf{cálculo diferencial}. Aquí $\Delta t$ actúa como un infinitésimo de tiempo: no es cero (porque no se puede dividir por cero), pero es \textit{arbitrariamente} cercano a cero.
  \end{itemize}
  \textbf{En resumen:} los límites nos dan una forma rigurosa de hablar de cantidades ``arbitrariamente pequeñas'' (infinitésimos), superando viejas paradojas como las de Zenón y abriendo la puerta a las dos grandes ramas del cálculo: \textbf{derivadas} (cómo cambian las cosas) e \textbf{integrales} (cómo se acumulan). Todo lo que practicarás en esta guía ---límites laterales, álgebra de límites, asíntotas--- son las herramientas para cruzar esa puerta con confianza.
\end{cuadroextra}

\vspace{8pt}

% ============================================================
\section{Introducción Intuitiva al Límite y Límites Laterales}
% ============================================================

\subsection{¿Qué es un límite?}

\begin{cuadroteoria}
  \textbf{Idea intuitiva:} El límite de $f(x)$ cuando $x$ se acerca a $a$, denotado $\lim_{x \to a} f(x)$, es el valor al que $f(x)$ se aproxima conforme $x$ se acerca a $a$ (por ambos lados), sin necesidad de que $f(a)$ esté definida.

  \textbf{Límites laterales:}
  \begin{itemize}[leftmargin=*]
    \item $\lim_{x \to a^-} f(x)$: límite por la \textbf{izquierda} (valores $x < a$).
    \item $\lim_{x \to a^+} f(x)$: límite por la \textbf{derecha} (valores $x > a$).
  \end{itemize}

  \textbf{Existencia del límite:}
  $$\boxed{\lim_{x \to a} f(x) = L \iff \lim_{x \to a^-} f(x) = L \;\text{ y }\; \lim_{x \to a^+} f(x) = L}$$
\end{cuadroteoria}

\begin{cuadrosolucion}
  \textbf{Observación importante:} El límite $\lim_{x \to a} f(x)$ \textbf{no depende} del valor de $f(a)$. La función puede estar definida o no en $x = a$, y el límite existe igual si los valores se aproximan a $L$.
\end{cuadrosolucion}

\subsection{Ejercicios resueltos}

\textbf{Ejemplo R1.} Evalúa $\lim_{x \to 2} (x^2 + 1)$ mediante una tabla de valores.

\begin{cuadrosolucion}
  \textbf{Solución:} Construimos una tabla con valores cercanos a 2:
  \begin{center}
    \small
    \begin{tabular}{cccccccc}
      \toprule
      $x$            & 1{,}9  & 1{,}99   & 1{,}999 & $\to 2^-$ & $\to 2^+$ & 2{,}001 & 2{,}01   \\
      \midrule
      $f(x) = x^2+1$ & 4{,}61 & 4{,}9601 & 4{,}996 & $\to 5$   &           & 5{,}004 & 5{,}0401 \\
      \bottomrule
    \end{tabular}
  \end{center}
  Los valores se acercan a 5 por ambos lados. Por tanto: $\lim_{x \to 2} (x^2 + 1) = 5$.
\end{cuadrosolucion}
\newpage
\textbf{Ejemplo R2.} Sea $f(x) = \dfrac{x^2 - 1}{x - 1}$ para $x \neq 1$. Evalúa $\lim_{x \to 1} f(x)$.

\begin{cuadrosolucion}
  \textbf{Solución:} Nota que $f(1)$ no está definida (división entre cero). Pero:
  $$\frac{x^2 - 1}{x - 1} = \frac{(x - 1)(x + 1)}{x - 1} = x + 1 \quad (x \neq 1).$$
  Entonces, cuando $x$ se acerca a 1, $f(x) = x + 1$ se acerca a 2. Por tanto:
  $$\lim_{x \to 1} \frac{x^2 - 1}{x - 1} = 2.$$
  Este es un ejemplo de \textbf{discontinuidad evitable}: la función tiene un hueco en $(1, 2)$, pero el límite existe.
\end{cuadrosolucion}

\textbf{Ejemplo R3.} Analiza la existencia del límite para $f(x) = |x|/x$ cuando $x \to 0$.

\begin{cuadrosolucion}
  \textbf{Solución:} Para $x > 0$: $\dfrac{|x|}{x} = 1$. Para $x < 0$: $\dfrac{|x|}{x} = -1$.
  \begin{itemize}
    \item $\lim_{x \to 0^+} \dfrac{|x|}{x} = 1$.
    \item $\lim_{x \to 0^-} \dfrac{|x|}{x} = -1$.
  \end{itemize}
  Como los límites laterales son distintos, $\lim_{x \to 0} \dfrac{|x|}{x}$ \textbf{no existe}.
\end{cuadrosolucion}

\textbf{Ejemplo R4.} Sea $f(x) = \begin{cases} x^2 + 1 & \text{si } x \ge 1 \\ 3x & \text{si } x < 1 \end{cases}$. ¿Existe $\lim_{x \to 1} f(x)$?

\begin{cuadrosolucion}
  \textbf{Solución:}
  \begin{itemize}
    \item Por la izquierda: $\lim_{x \to 1^-} 3x = 3$.
    \item Por la derecha: $\lim_{x \to 1^+} (x^2 + 1) = 2$.
  \end{itemize}
  Son distintos, por lo que el límite no existe.
\end{cuadrosolucion}

\newpage

\subsection{Ejercicios propuestos}

\begin{enumerate}[leftmargin=*]
  \item \facil{} Evalúa mediante tabla de valores: $\lim_{x \to 3} (2x + 1)$.

  \item \facil{} Evalúa mediante tabla de valores: $\lim_{x \to 0} \cos x$.

  \item \facil{} Evalúa: $\lim_{x \to 2} \dfrac{x^2 - 4}{x - 2}$ (pista: factoriza).

  \item \facil{} Determina si el límite existe: $f(x) = \begin{cases} x + 2 & \text{si } x < 1 \\ 4 & \text{si } x = 1 \\ 2x + 1 & \text{si } x > 1 \end{cases}$, en $x = 1$.

  \item \medio{} Determina si existe $\lim_{x \to 0} \dfrac{1}{x^2}$.

  \item \medio{} Determina si existe $\lim_{x \to 2} \dfrac{2x - 4}{|x - 2|}$.

  \item \medio{} Sea $f(x) = \begin{cases} x^2 & \text{si } x < 0 \\ 1 & \text{si } x = 0 \\ x + 2 & \text{si } x > 0 \end{cases}$. Determina si existe $\lim_{x \to 0} f(x)$.

  \item \medio{} Determina si existe $\lim_{x \to 4} \dfrac{x - 4}{\sqrt{x} - 2}$ (pista: racionaliza).

  \item \dificil{} Dibuja una función que cumpla simultáneamente: $f(2) = 5$, $\lim_{x \to 2^-} f(x) = 3$, $\lim_{x \to 2^+} f(x) = 3$ y $\lim_{x \to 3} f(x) = \infty$.

  \item \dificil{} Encuentra una fórmula para una función que tenga un hueco en $x = -1$, salto en $x = 0$ y esté definida en todos los reales.

  \item \dificil{} Sea $f(x) = \begin{cases} 3x - a & \text{si } x < 2 \\ ax + 1 & \text{si } x \ge 2 \end{cases}$. Encuentra $a$ para que $\lim_{x \to 2} f(x)$ exista.

  \item \dificil{} Analiza la existencia de $\lim_{x \to 1} \dfrac{|x - 1|}{x - 1}$ y justifica tu respuesta.
\end{enumerate}

\newpage

% ============================================================
\section{Cálculo de Límites Algebraicos}
% ============================================================

\subsection{Propiedades de los límites}

\begin{cuadroteoria}
  Si $\lim_{x \to a} f(x) = L$ y $\lim_{x \to a} g(x) = M$, entonces:
  \begin{align*}
    \lim_{x \to a} [f(x) + g(x)]                  & = L + M                                                               \\
    \lim_{x \to a} [f(x) - g(x)]                  & = L - M                                                               \\
    \lim_{x \to a} [f(x) \cdot g(x)]              & = L \cdot M                                                           \\
    \lim_{x \to a} \left[\frac{f(x)}{g(x)}\right] & = \frac{L}{M} \quad (M \neq 0)                                        \\
    \lim_{x \to a} [c \cdot f(x)]                 & = c \cdot L                                                           \\
    \lim_{x \to a} [f(x)]^n                       & = L^n                                                                 \\
    \lim_{x \to a} \sqrt[n]{f(x)}                 & = \sqrt[n]{L} \quad (\text{si } L \geq 0 \text{ para } n \text{ par})
  \end{align*}

  \textbf{Límite de una función polinomial:}
  $$\lim_{x \to a} P(x) = P(a)$$

  \textbf{Sustitución directa:} para funciones polinomiales y racionales (cuando el denominador no se anula), el límite se evalúa sustituyendo $x = a$.
\end{cuadroteoria}

\begin{cuadroadvertencia}
  \textbf{Forma indeterminada $0/0$:} Si al sustituir obtienes $\frac{0}{0}$, necesitas \textbf{simplificar} antes de evaluar. Técnicas:
  \begin{enumerate}[leftmargin=*]
    \item \textbf{Factorizar} el numerador y/o denominador, y cancelar el factor común.
    \item \textbf{Racionalizar} (multiplicar por el conjugado) cuando hay raíces.
    \item En límites de cocientes de polinomios, el factor que causa el problema siempre se cancela o se llega a un límite infinito.
  \end{enumerate}
\end{cuadroadvertencia}

\subsection{Ejercicios resueltos}

\textbf{Ejemplo R1.} Calcula $\lim_{x \to 3} (2x^2 - 5x + 1)$.

\begin{cuadrosolucion}
  \textbf{Solución:} Sustitución directa:
  $$2(3)^2 - 5(3) + 1 = 18 - 15 + 1 = 4.$$
\end{cuadrosolucion}

\textbf{Ejemplo R2.} Calcula $\lim_{x \to 2} \dfrac{x^2 - 5x + 6}{x - 2}$.

\begin{cuadrosolucion}
  \textbf{Solución:} Sustituyendo $x = 2$ llegamos a $0/0$. Factorizamos:
  $$\frac{x^2 - 5x + 6}{x - 2} = \frac{(x - 2)(x - 3)}{x - 2} = x - 3 \quad (x \neq 2).$$
  Entonces: $\lim_{x \to 2} \dfrac{x^2 - 5x + 6}{x - 2} = \lim_{x \to 2} (x - 3) = -1$.
\end{cuadrosolucion}

\textbf{Ejemplo R3.} Calcula $\lim_{x \to 4} \dfrac{\sqrt{x} - 2}{x - 4}$.

\begin{cuadrosolucion}
  \textbf{Solución:} Sustituyendo da $0/0$. Racionalizamos el numerador multiplicando por el conjugado $\sqrt{x} + 2$:
  $$\frac{(\sqrt{x} - 2)(\sqrt{x} + 2)}{(x - 4)(\sqrt{x} + 2)} = \frac{x - 4}{(x - 4)(\sqrt{x} + 2)} = \frac{1}{\sqrt{x} + 2}.$$
  Entonces:
  $$\lim_{x \to 4} \dfrac{\sqrt{x} - 2}{x - 4} = \lim_{x \to 4} \dfrac{1}{\sqrt{x} + 2} = \dfrac{1}{\sqrt{4} + 2} = \dfrac{1}{4}.$$
\end{cuadrosolucion}

\textbf{Ejemplo R4.} Calcula $\lim_{x \to 1} \dfrac{x^3 - 1}{x - 1}$.

\begin{cuadrosolucion}
  \textbf{Solución:} Factorizamos la diferencia de cubos $x^3 - 1 = (x - 1)(x^2 + x + 1)$:
  $$\frac{x^3 - 1}{x - 1} = x^2 + x + 1 \quad (x \neq 1).$$
  $$\lim_{x \to 1} \dfrac{x^3 - 1}{x - 1} = \lim_{x \to 1} (x^2 + x + 1) = 1 + 1 + 1 = 3.$$
\end{cuadrosolucion}

\textbf{Ejemplo R5.} Calcula $\lim_{x \to 0} \dfrac{\sqrt{1 + x} - 1}{x}$.

\begin{cuadrosolucion}
  \textbf{Solución:} Racionalizamos:
  $$\frac{(\sqrt{1 + x} - 1)(\sqrt{1 + x} + 1)}{x(\sqrt{1 + x} + 1)} = \frac{1 + x - 1}{x(\sqrt{1 + x} + 1)} = \frac{x}{x(\sqrt{1 + x} + 1)} = \frac{1}{\sqrt{1 + x} + 1}.$$
  $$\lim_{x \to 0} \dfrac{\sqrt{1 + x} - 1}{x} = \dfrac{1}{\sqrt{1} + 1} = \dfrac{1}{2}.$$
  Este límite aparecerá de nuevo en la definición de derivada.
\end{cuadrosolucion}

\newpage

\subsection{Ejercicios propuestos}

\begin{enumerate}[leftmargin=*, resume]
  \item \facil{} Calcula: $\lim_{x \to 1} (x^3 - 2x^2 + x)$.

  \item \facil{} Calcula: $\lim_{x \to -2} \dfrac{x^2 + x - 2}{x + 2}$.

  \item \facil{} Calcula: $\lim_{x \to 5} \dfrac{x^2 - 25}{x - 5}$.

  \item \facil{} Calcula: $\lim_{x \to 3} \dfrac{2x - 6}{x^2 - 9}$.

  \item \medio{} Calcula: $\lim_{x \to 9} \dfrac{\sqrt{x} - 3}{x - 9}$.

  \item \medio{} Calcula: $\lim_{x \to 0} \dfrac{\sqrt{4 + x} - 2}{x}$.

  \item \medio{} Calcula: $\lim_{x \to 2} \dfrac{x^3 - 8}{x - 2}$.

  \item \medio{} Calcula: $\lim_{x \to 1} \dfrac{x^2 + 2x - 3}{x^2 - 1}$.

  \item \dificil{} Calcula: $\lim_{x \to 0} \dfrac{1 - \sqrt{1 - x}}{x}$.

  \item \dificil{} Calcula: $\lim_{x \to 1} \dfrac{\sqrt{x} - 1}{x - 1}$ (pista: haz $t = \sqrt{x}$, con $t \to 1$).

  \item \dificil{} Calcula: $\lim_{h \to 0} \dfrac{(2 + h)^3 - 8}{h}$.

  \item \dificil{} Calcula: $\lim_{x \to 0} \dfrac{3x - \sin(3x)}{x}$ (pista: usa el límite fundamental que verás en la siguiente sección).
\end{enumerate}

\newpage

% ============================================================
\section{Límites Trigonométricos y Límites al Infinito}
% ============================================================

\subsection{Límite fundamental y límites al infinito}

\begin{cuadroteoria}
  \textbf{Límite trigonométrico fundamental:}
  $$\boxed{\lim_{x \to 0} \frac{\sin x}{x} = 1}$$

  \textbf{Variantes útiles:}
  \begin{itemize}[leftmargin=*]
    \item $\lim_{x \to 0} \dfrac{\sin(kx)}{kx} = 1$ (para cualquier constante $k \neq 0$).
    \item $\lim_{x \to 0} \dfrac{1 - \cos x}{x} = 0$.
    \item $\lim_{x \to 0} \dfrac{\tan x}{x} = 1$.
  \end{itemize}
\end{cuadroteoria}

\begin{cuadrosolucion}
  \textbf{Límites al infinito:}
  $$\lim_{x \to \infty} f(x) = L$$
  significa que $f(x)$ se acerca a $L$ conforme $x$ crece sin límite.

  \textbf{Técnica para funciones racionales en el infinito:} divide numerador y denominador por la mayor potencia de $x$ en el denominador.

  \textbf{Asíntotas horizontales:} si $\lim_{x \to \infty} f(x) = L$, entonces $y = L$ es asíntota horizontal.
\end{cuadrosolucion}

\begin{cuadroadvertencia}
  \textbf{Límites infinitos:} $\lim_{x \to a} f(x) = \infty$ o $-\infty$ indica que la función crece o decrece sin límite cerca de $x = a$. Representan asíntotas verticales.
\end{cuadroadvertencia}

\begin{cuadroextra}
  \textbf{Teorema del Emparedado} (\textit{Squeeze Theorem})\textbf{:} Si $g(x) \leq f(x) \leq h(x)$ para todo $x$ cerca de $a$ (excepto posiblemente en $a$), y
  $$\lim_{x \to a} g(x) = \lim_{x \to a} h(x) = L,$$
  entonces $\lim_{x \to a} f(x) = L$.

  Este teorema se utiliza para demostrar el límite fundamental $\lim_{x \to 0} \frac{\sin x}{x} = 1$, encuadrando $\frac{\sin x}{x}$ entre funciones cuyo límite es 1.
\end{cuadroextra}

\subsection{Ejercicios resueltos}

\textbf{Ejemplo R1.} Calcula $\lim_{x \to 0} \dfrac{\sin(3x)}{x}$.

\begin{cuadrosolucion}
  \textbf{Solución:} Multiplicamos y dividimos por 3:
  $$\frac{\sin(3x)}{x} = 3 \cdot \frac{\sin(3x)}{3x}.$$
  Como $\lim_{x \to 0} \frac{\sin(3x)}{3x} = 1$ (variante del límite fundamental con $k = 3$):
  $$\lim_{x \to 0} \frac{\sin(3x)}{x} = 3.$$
\end{cuadrosolucion}

\textbf{Ejemplo R2.} Calcula $\lim_{x \to 0} \dfrac{\tan x}{x}$.

\begin{cuadrosolucion}
  \textbf{Solución:} Escribimos $\tan x = \dfrac{\sin x}{\cos x}$:
  $$\frac{\tan x}{x} = \frac{\sin x}{x} \cdot \frac{1}{\cos x}.$$
  $$\lim_{x \to 0} \frac{\tan x}{x} = 1 \cdot \frac{1}{\cos 0} = 1.$$
\end{cuadrosolucion}

\textbf{Ejemplo R3.} Calcula $\lim_{x \to \infty} \dfrac{3x^2 + 2x - 1}{5x^2 + x + 4}$.

\begin{cuadrosolucion}
  \textbf{Solución:} Dividimos por $x^2$:
  $$\frac{3 + \frac{2}{x} - \frac{1}{x^2}}{5 + \frac{1}{x} + \frac{4}{x^2}} \xrightarrow{x \to \infty} \frac{3 + 0 - 0}{5 + 0 + 0} = \frac{3}{5}.$$
\end{cuadrosolucion}

\textbf{Ejemplo R4.} Calcula $\lim_{x \to \infty} \dfrac{2x^2 + 3}{5x^3 - x}$.

\begin{cuadrosolucion}
  \textbf{Solución:} Dividimos por $x^3$:
  $$\frac{\frac{2}{x} + \frac{3}{x^3}}{5 - \frac{1}{x^2}} \xrightarrow{x \to \infty} \frac{0 + 0}{5 - 0} = 0.$$
  Asíntota horizontal: $y = 0$.
\end{cuadrosolucion}
\newpage
\textbf{Ejemplo R5.} Calcula $\lim_{x \to \infty} \dfrac{3x^3 - 2x + 1}{x^2 + 1}$.

\begin{cuadrosolucion}
  \textbf{Solución:} Dividimos por $x^3$ (la mayor potencia del denominador es $x^2$, así que también funciona observar que el numerador crece más rápido):

  Método con división por $x^2$:
  $$\frac{3x - \frac{2}{x} + \frac{1}{x^2}}{1 + \frac{1}{x^2}} \xrightarrow{x \to \infty} \frac{3x - 0 + 0}{1 + 0} \to +\infty.$$
  El límite es $\infty$ (la función crece sin límite).
\end{cuadrosolucion}

\textbf{Ejemplo R6.} Calcula $\lim_{x \to 3} \dfrac{2}{x - 3}$ (límite infinito).

\begin{cuadrosolucion}
  \textbf{Solución:} Cuando $x$ se acerca a 3 por la derecha, $x - 3 \to 0^+$ y $\frac{2}{x-3} \to +\infty$. Cuando $x$ se acerca por la izquierda, $x - 3 \to 0^-$ y $\frac{2}{x-3} \to -\infty$. Por tanto, el límite bilateral no existe, pero hay asíntota vertical en $x = 3$:
  $$\lim_{x \to 3^+} \frac{2}{x - 3} = +\infty, \qquad \lim_{x \to 3^-} \frac{2}{x - 3} = -\infty.$$
\end{cuadrosolucion}

\newpage

\subsection{Ejercicios propuestos}

\begin{enumerate}[leftmargin=*, resume]
  \item \facil{} Calcula: $\lim_{x \to 0} \dfrac{\sin(2x)}{2x}$

  \item \facil{} Calcula: $\lim_{x \to 0} \dfrac{\sin(5x)}{x}$

  \item \facil{} Calcula: $\lim_{x \to 0} \dfrac{1 - \cos x}{x}$

  \item \facil{} Calcula: $\lim_{x \to \infty} \dfrac{1}{x}$

  \item \medio{} Calcula: $\lim_{x \to 0} \dfrac{\sin(4x)}{\sin(2x)}$

  \item \medio{} Calcula: $\lim_{x \to 0} \dfrac{1 - \cos x}{x^2}$ (pista: multiplica por el conjugado $1 + \cos x$)

  \item \medio{} Calcula: $\lim_{x \to \infty} \dfrac{7x^2 - 3}{2x^2 + 1}$

  \item \medio{} Calcula: $\lim_{x \to \infty} \dfrac{x + 5}{x^2 + 1}$

  \item \dificil{} Calcula: $\lim_{x \to 0} \dfrac{\sin(3x)}{\tan(2x)}$

  \item \dificil{} Calcula: $\lim_{x \to \infty} \dfrac{4x^3 + 2}{5x^3 - x^2 + 1}$

  \item \dificil{} Calcula: $\lim_{x \to \infty} \dfrac{x^2 + 4x}{x + 1}$ ¿Existe asíntota horizontal u oblicua?

  \item \dificil{} Calcula: $\lim_{x \to 0} \dfrac{\sin^2 x}{x^2}$
\end{enumerate}

\newpage

% ============================================================
\section{Continuidad}
% ============================================================

\subsection{Definición y tipos de discontinuidades}

\begin{cuadroteoria}
  \textbf{Continuidad en un punto:} $f$ es continua en $x = a$ si se cumplen tres condiciones:
  \begin{enumerate}
    \item $f(a)$ está definida.
    \item $\lim_{x \to a} f(x)$ existe.
    \item $\lim_{x \to a} f(x) = f(a)$.
  \end{enumerate}
\end{cuadroteoria}

\begin{cuadrosolucion}
  \textbf{Tipos de discontinuidad:}
  \begin{itemize}[leftmargin=*]
    \item \textbf{Evitable (hueco):} $\lim_{x \to a} f(x)$ existe pero no es igual a $f(a)$, o $f(a)$ no está definida. Se puede ``reparar'' definiendo $f(a)$ apropiadamente.
    \item \textbf{De salto:} los límites laterales existen pero son distintos.
    \item \textbf{Infinita:} al menos uno de los límites laterales es $\pm\infty$ (asíntota vertical).
  \end{itemize}
\end{cuadrosolucion}

\begin{cuadroadvertencia}
  \textbf{Continuidad en un intervalo:} $f$ es continua en $(a, b)$ si es continua en cada punto del intervalo. Las funciones polinomiales son continuas en $\mathbb{R}$; las racionales, en todo su dominio; $\sin x$, $\cos x$, $e^x$ y $\ln x$ (en su dominio) son continuas.
\end{cuadroadvertencia}

\begin{cuadroextra}
  \textbf{Teorema del Valor Intermedio (TVI):} Si $f$ es continua en $[a, b]$ y $N$ es cualquier número entre $f(a)$ y $f(b)$, entonces existe al menos un $c \in (a, b)$ tal que $f(c) = N$.

  \textbf{Consecuencia para raíces:} Si $f$ es continua en $[a, b]$ y $f(a)$ y $f(b)$ tienen \textbf{signos opuestos}, entonces $f$ tiene al menos una raíz en $(a, b)$.
\end{cuadroextra}

\subsection{Ejercicios resueltos}

\textbf{Ejemplo R1.} Determina si $f(x) = \dfrac{x^2 - 9}{x - 3}$ (para $x \neq 3$) es continua en $x = 3$.

\begin{cuadrosolucion}
  \textbf{Solución:} $f(3)$ no está definida, por lo que falla la primera condición. Hay una discontinuidad \textbf{evitable}. Si definimos $f(3) = 6$, la función se vuelve continua (pues $\lim_{x \to 3} f(x) = 6$).
\end{cuadrosolucion}
\newpage
\textbf{Ejemplo R2.} ¿En qué puntos es discontinua $f(x) = \dfrac{x + 2}{x^2 - 4}$?

\begin{cuadrosolucion}
  \textbf{Solución:} El denominador se anula en $x = \pm 2$.
  \begin{itemize}
    \item En $x = 2$: el numerador no se anula. Hay asíntota vertical (discontinuidad infinita).
    \item En $x = -2$: $x^2 - 4 = (x - 2)(x + 2)$, el numerador $x + 2$ también se anula. La función se simplifica: $f(x) = \dfrac{1}{x - 2}$ para $x \neq -2$. El límite en $x = -2$ es $\dfrac{1}{-4} = -\dfrac{1}{4}$, pero $f(-2)$ no está definida: discontinuidad evitable.
  \end{itemize}
\end{cuadrosolucion}

\textbf{Ejemplo R3.} Encuentra el valor de $k$ para que $f$ sea continua en $x = 1$:
$$f(x) = \begin{cases} 2x + k & \text{si } x < 1 \\ 5x - 1 & \text{si } x \ge 1 \end{cases}$$

\begin{cuadrosolucion}
  \textbf{Solución:} Para que sea continua, los límites laterales deben coincidir y ser igual a $f(1)$:
  \begin{itemize}
    \item Por la izquierda: $\lim_{x \to 1^-} (2x + k) = 2 + k$.
    \item Por la derecha: $f(1) = 5(1) - 1 = 4$.
  \end{itemize}
  Igualamos: $2 + k = 4 \Rightarrow k = 2$.
\end{cuadrosolucion}

\textbf{Ejemplo R4.} Determina los intervalos de continuidad de $f(x) = \dfrac{1}{x^2 - 1}$.

\begin{cuadrosolucion}
  \textbf{Solución:} El denominador se anula cuando $x^2 - 1 = 0 \Rightarrow x = \pm 1$. La función es continua en todo su dominio:
  $$(-\infty, -1) \cup (-1, 1) \cup (1, \infty).$$
  Tiene asíntotas verticales en $x = -1$ y $x = 1$.
\end{cuadrosolucion}

\newpage

\subsection{Ejercicios propuestos}

\begin{enumerate}[leftmargin=*, resume]
  \item \facil{} Determina si $f(x) = x^2 - 3x + 2$ es continua en $x = 2$ (justifica las 3 condiciones).

  \item \facil{} Determina los puntos de discontinuidad de $f(x) = \dfrac{x + 1}{x - 3}$.

  \item \facil{} Determina si la siguiente función es continua en $x = 0$:
        $$f(x) = \begin{cases} x^2 + 1 & \text{si } x < 0 \\ 2 & \text{si } x = 0 \\ x + 1 & \text{si } x > 0 \end{cases}$$

  \item \facil{} Determina los intervalos de continuidad de $f(x) = \sqrt{x - 2}$.

  \item \medio{} Encuentra $k$ para que $f$ sea continua en $x = 2$:
        $$f(x) = \begin{cases} kx + 3 & \text{si } x \le 2 \\ x^2 - 1 & \text{si } x > 2 \end{cases}$$

  \item \medio{} Determina el tipo de discontinuidad de $f(x) = \dfrac{x^2 - 16}{x + 4}$ en $x = -4$, y repara la función si es posible.

  \item \medio{} Determina los intervalos de continuidad de $f(x) = \dfrac{x}{x^2 + x - 6}$.

  \item \medio{} Sea $f(x) = \begin{cases} \dfrac{\sin x}{x} & \text{si } x \neq 0 \\ 1 & \text{si } x = 0 \end{cases}$. ¿Es continua en $x = 0$?

  \item \dificil{} Encuentra $a$ y $b$ para que $f$ sea continua en todo $\mathbb{R}$:
        $$f(x) = \begin{cases} ax^2 + b & \text{si } x < 1 \\ 5x - 2 & \text{si } 1 \le x \le 3 \\ 2a x + b & \text{si } x > 3 \end{cases}$$

  \item \dificil{} Usa el Teorema del Valor Intermedio para mostrar que $f(x) = x^3 - 4x + 1$ tiene al menos una raíz en el intervalo $[0, 1]$.

  \item \dificil{} Determina si $f(x) = \begin{cases} \dfrac{x^2 - 4}{x - 2} & \text{si } x \neq 2 \\ 3 & \text{si } x = 2 \end{cases}$ es continua en $x = 2$. Si no lo es, ¿puedes repararla?

  \item \dificil{} Demuestra que la ecuación $\cos x = x$ tiene al menos una solución en $[0, 1]$.
\end{enumerate}

\newpage

% ============================================================
\section{Clave de Respuestas}
% ============================================================

\subsection*{Sección 1: Introducción (Ejercicios 1--12)}

\begin{enumerate}[leftmargin=*, label=\textbf{\arabic*.}]
  \item $\lim_{x \to 3} (2x + 1) = 7$.

  \item $\lim_{x \to 0} \cos x = 1$.

  \item $\dfrac{x^2 - 4}{x - 2} = x + 2$, entonces el límite es $4$.

  \item Por izquierda: $3$; por derecha (y $f(1)$): $3$. Los límites laterales coinciden: el límite existe y vale 3.

  \item $x^2 \to 0^+$ siempre (es cuadrado), así que $\dfrac{1}{x^2} \to +\infty$. El límite es $+\infty$; no existe como límite finito.

  \item Por derecha: $\dfrac{2(x-2)}{x-2} = 2 \to 2$. Por izquierda: $\dfrac{2(x-2)}{-(x-2)} = -2 \to -2$. No existe (laterales distintos).

  \item Por izquierda: $\lim_{x \to 0^-} x^2 = 0$. Por derecha: $\lim_{x \to 0^+} (x + 2) = 2$. Laterales distintos, no existe.

  \item Racionalizamos: $\dfrac{(x - 4)(\sqrt{x} + 2)}{(\sqrt{x} - 2)(\sqrt{x} + 2)} = \dfrac{(x-4)(\sqrt{x}+2)}{x - 4} = \sqrt{x} + 2$. Límite: $\sqrt{4} + 2 = 4$.

  \item (Gráfico, respuesta abierta). Una función con hueco en $(2,3)$, salto en algún punto antes de 3, y asíntota vertical en $x = 3$.

  \item Ejemplo: $f(x) = \begin{cases} x + 2 & \text{si } x \neq -1, 0 \\ 1 & \text{si } x = 0 \end{cases}$ (hueco en $x = -1$: no definida; salto en $x = 0$). Otras respuestas son válidas.

  \item $\lim_{x \to 2^-} (3x - a) = 6 - a$; $\lim_{x \to 2^+} (ax + 1) = 2a + 1$. Igualando: $6 - a = 2a + 1 \Rightarrow 3a = 5 \Rightarrow a = \dfrac{5}{3}$.

  \item Para $x > 1$: $\dfrac{x-1}{x-1} = 1$; para $x < 1$: $\dfrac{-(x-1)}{x-1} = -1$. Laterales distintos, el límite no existe.
\end{enumerate}

\subsection*{Sección 2: Límites Algebraicos (Ejercicios 13--24)}

\begin{enumerate}[leftmargin=*, label=\textbf{\arabic*.}]
  \setcounter{enumi}{12}
  \item $1 - 2 + 1 = 0$.

  \item $\dfrac{x^2 + x - 2}{x + 2} = \dfrac{(x + 2)(x - 1)}{x + 2} = x - 1$. En $x = -2$: $-3$.

  \item $x + 5$; límite: $10$.

  \item $\dfrac{2(x - 3)}{(x - 3)(x + 3)} = \dfrac{2}{x + 3}$; en $x = 3$: $\dfrac{2}{6} = \dfrac{1}{3}$.

  \item $\dfrac{1}{\sqrt{x} + 3}$; en $x = 9$: $\dfrac{1}{6}$.

  \item $\dfrac{1}{\sqrt{4 + x} + 2}$; en $x = 0$: $\dfrac{1}{4}$.

  \item $x^2 + 2x + 4$; en $x = 2$: $12$.

  \item $\dfrac{(x - 1)(x + 3)}{(x - 1)(x + 1)} = \dfrac{x + 3}{x + 1}$; en $x = 1$: $\dfrac{4}{2} = 2$.

  \item Racionalizamos: $\dfrac{x}{(\sqrt{1 - x} + 1)(1 + \sqrt{1-x})} \dots$ Al final: $\lim_{x \to 0} \dfrac{1 - \sqrt{1 - x}}{x} = \dfrac{1}{2}$.

  \item Con $t = \sqrt{x}$: $x = t^2$, cuando $x \to 1$, $t \to 1$: $\dfrac{t - 1}{t^2 - 1} = \dfrac{1}{t + 1}$. Límite: $\dfrac{1}{2}$.

  \item Expandimos: $(2 + h)^3 = 8 + 12h + 6h^2 + h^3$. Restando 8 y dividiendo entre $h$: $12 + 6h + h^2$. Límite: $12$.

  \item $\dfrac{3x - \sin(3x)}{x} = 3 - \dfrac{\sin(3x)}{x} = 3 - 3\left(\dfrac{\sin(3x)}{3x}\right) \to 3 - 3 = 0$.
\end{enumerate}

\subsection*{Sección 3: Trigonométricos e Infinito (Ejercicios 25--36)}

\begin{enumerate}[leftmargin=*, label=\textbf{\arabic*.}]
  \setcounter{enumi}{24}
  \item $1$ (es el límite fundamental con $k = 2$).

  \item $5 \cdot \lim_{x \to 0} \dfrac{\sin(5x)}{5x} = 5$.

  \item $0$.

  \item $0$.

  \item $\dfrac{\sin(4x)}{\sin(2x)} = \dfrac{4 \cdot \frac{\sin(4x)}{4x}}{2 \cdot \frac{\sin(2x)}{2x}} \to \dfrac{4}{2} = 2$.

  \item $\dfrac{1 - \cos x}{x^2} \cdot \dfrac{1 + \cos x}{1 + \cos x} = \dfrac{\sin^2 x}{x^2(1 + \cos x)} = \left(\dfrac{\sin x}{x}\right)^2 \cdot \dfrac{1}{1 + \cos x} \to 1 \cdot \dfrac{1}{2} = \dfrac{1}{2}$.

  \item $\dfrac{7}{2}$.

  \item $0$.

  \item $\dfrac{\sin(3x)}{\tan(2x)} = \dfrac{\sin(3x)}{\sin(2x)/\cos(2x)} = \dfrac{\sin(3x) \cos(2x)}{\sin(2x)}$. Dividiendo todo por $x$: $\dfrac{3 \cdot \frac{\sin(3x)}{3x} \cdot \cos(2x)}{2 \cdot \frac{\sin(2x)}{2x}} \to \dfrac{3 \cdot 1 \cdot 1}{2 \cdot 1} = \dfrac{3}{2}$.

  \item $\dfrac{4 + \frac{2}{x^3}}{5 - \frac{1}{x^2} + \frac{1}{x^3}} \to \dfrac{4}{5}$.

  \item Dividimos por $x^2$: $\dfrac{1 + \frac{4}{x}}{\frac{1}{x} + \frac{1}{x^2}} \to \infty$. No hay asíntota horizontal; hay asíntota oblicua porque la diferencia de grados es 1. Dividiendo: $\dfrac{x^2 + 4x}{x + 1} = x + 3 - \dfrac{3}{x + 1}$; asíntota oblicua: $y = x + 3$.

  \item $\left(\dfrac{\sin x}{x}\right)^2 \to 1^2 = 1$.
\end{enumerate}

\subsection*{Sección 4: Continuidad (Ejercicios 37--48)}

\begin{enumerate}[leftmargin=*, label=\textbf{\arabic*.}]
  \setcounter{enumi}{36}
  \item Sí: $f(2) = 4 - 6 + 2 = 0$; $\lim_{x \to 2} f(x) = 0$; coinciden.

  \item Discontinua en $x = 3$ (asíntota vertical). Continua en $\mathbb{R} \setminus \{3\}$.

  \item Por izquierda: $1$; por derecha y en $x = 0$: $1$. Coinciden, sí es continua en $x = 0$.

  \item $[2, \infty)$.

  \item $\lim_{x \to 2^-} (kx + 3) = 2k + 3$; $f(2) = 2k + 3$. $\lim_{x \to 2^+} (x^2 - 1) = 3$. Igualamos: $2k + 3 = 3 \Rightarrow k = 0$.

  \item En $x = -4$: $f(x) = \dfrac{(x - 4)(x + 4)}{x + 4} = x - 4$ para $x \neq -4$. El límite es $-8$, pero $f(-4)$ no existe. Discontinuidad evitable; se repara definiendo $f(-4) = -8$.

  \item $x^2 + x - 6 = (x + 3)(x - 2)$. Dominio: $(-\infty, -3) \cup (-3, 2) \cup (2, \infty)$. La función es continua en todo su dominio.

  \item $\lim_{x \to 0} \dfrac{\sin x}{x} = 1 = f(0)$. Sí es continua.

  \item En $x = 1$: $a + b = 3$. En $x = 3$: $13 = 6a + b$. Restando: $5a = 10 \Rightarrow a = 2$; $b = 1$.

  \item $f(0) = 1 > 0$; $f(1) = 1 - 4 + 1 = -2 < 0$. Por el TVI, existe $c \in (0, 1)$ con $f(c) = 0$.

  \item La función simplificada es $f(x) = x + 2$ para $x \neq 2$. El límite en $x = 2$ es 4 y $f(2) = 3$, así que NO es continua. Se puede reparar cambiando $f(2) = 4$.

  \item Sea $f(x) = \cos x - x$. $f(0) = 1 > 0$; $f(1) \approx 0{,}54 - 1 = -0{,}46 < 0$. Por TVI, existe $c \in (0, 1)$ con $f(c) = 0$, es decir, $\cos c = c$.
\end{enumerate}

\end{document}